\chapter{The low individual degree test}
\label{chap:test}
\label{sec:the-test}

\section{The game and its strategies}

\begin{definition}[Roles]\label{def:roles}
	\lean{MIPStarRE.LDT.Role, MIPStarRE.LDT.Role.other}
	\leanok
	A role is an element of $\{\mathrm A,\mathrm B\}$. For a role $r \in \{\mathrm A,\mathrm B\}$, write $\overline r$ for the other element of $\{\mathrm A,\mathrm B\}$.
\end{definition}

\begin{definition}[Low individual degree test]\label{def:low-individual-degree-test}
	\lean{MIPStarRE.LDT.ProjStrat.lowIndividualDegreeFailureProbability, MIPStarRE.LDT.ProjStrat.PassesLowIndividualDegreeTest}
	\leanok
	\uses{def:roles}
	Fix integers $m,d \ge 0$ and a prime power $q$. The $(m,q,d)$ low individual degree test has three equally likely subtests.
	In the axis-parallel lines test, choose a uniformly random role $r \in \{\mathrm A,\mathrm B\}$, a uniformly random point $u \in \F_q^m$, and a uniformly random coordinate $i \in \{1,\dots,m\}$. Let $\ell=\{u+t e_i : t \in \F_q\}$. Player $r$ receives $\ell$ and returns a degree-$d$ univariate polynomial $f\colon \ell \to \F_q$; Player $\overline r$ receives $u$ and returns a field element $a \in \F_q$. The verifier accepts when $f(u)=a$.
	In the self-consistency test, both provers receive the same uniformly random point $u \in \F_q^m$ and must return the same field element.
	In the diagonal lines test, choose a uniformly random role $r \in \{\mathrm A,\mathrm B\}$, a uniformly random point $u \in \F_q^m$, a uniformly random index $i \in \{1,\dots,m\}$, and a uniformly random direction vector $v \in \F_q^m$ whose last $m-i$ coordinates are $0$. Let $\ell=\{u+t v : t \in \F_q\}$. Player $r$ receives $\ell$ and returns a degree-$md$ univariate polynomial $f\colon \ell \to \F_q$; Player $\overline r$ receives $u$ and returns a field element $a \in \F_q$. The verifier accepts when $f(u)=a$.
\end{definition}

\begin{lemma}[Branch-average form of the classical test]\label{fig:test}
	\lean{MIPStarRE.LDT.Test.TwoProverClassicalLIDStrategy.lowIndividualDegreeAcceptanceProbability_eq_branchAverage}
	\leanok
	\uses{def:low-individual-degree-test}
	For every classical strategy, the acceptance probability of the low individual degree test is
	\[
		\frac13\left(\frac{P^{\mathrm{line,point}}_{\mathrm{axis}}+P^{\mathrm{point,line}}_{\mathrm{axis}}}{2}
		+P_{\mathrm{self}}
		+\frac{P^{\mathrm{line,point}}_{\mathrm{diag}}+P^{\mathrm{point,line}}_{\mathrm{diag}}}{2}\right),
	\]
	where the five terms are the acceptance probabilities of the two axis-parallel role orderings, the self-consistency branch, and the two diagonal role orderings.
\end{lemma}
\begin{proof}
	\leanok
	\proves{fig:test}
	\uses{def:low-individual-degree-test}
	Expand the definition of the total acceptance probability as the average of
	the three subtests.  The axis-parallel and diagonal branches each first
	choose one of the two role orderings uniformly, so their branch acceptance
	probabilities are the corresponding averages displayed above.
\end{proof}

\begin{definition}[Symmetric projective strategy]\label{def:symmetric-projective-strategy}
	\lean{MIPStarRE.LDT.SymStrat}
	\leanok
	A symmetric projective strategy for the $(m,q,d)$ low individual degree test consists of a permutation-invariant bipartite state $\ket{\psi} \in \mathcal H \ot \mathcal H$, a projective point measurement $A^u=\{A^u_a\}_{a \in \F_q}$ for each $u \in \F_q^m$, a projective axis-parallel line measurement $B^\ell=\{B^\ell_f\}$ for each axis-parallel line $\ell$, and a projective diagonal-line measurement $L^\ell=\{L^\ell_f\}$ for each line $\ell$.
\end{definition}

\begin{definition}[General projective strategy]\label{def:projective-strategy}
	\lean{MIPStarRE.LDT.ProjStrat}
	\leanok
	A general projective strategy for the $(m,q,d)$ low individual degree test consists of a bipartite state $\ket{\psi} \in \mathcal H_{\mathrm A} \ot \mathcal H_{\mathrm B}$ together with point, axis-parallel line, and diagonal-line projective measurements for each prover separately. We write these measurements as $A^{\mathrm A}, B^{\mathrm A}, L^{\mathrm A}$ on the first prover and $A^{\mathrm B}, B^{\mathrm B}, L^{\mathrm B}$ on the second.
\end{definition}

\begin{lemma}[Bipartite strategy block identities]\label{lem:bipartite-strategy-block-identities}
	\lean{MIPStarRE.LDT.ProjStrat.localDirectSumBlock_nonneg, MIPStarRE.LDT.ProjStrat.roleBlock_mul, MIPStarRE.LDT.ProjStrat.roleBlock_nonneg}
	\leanok
	\uses{def:projective-strategy}
	The block operators obtained from a bipartite projective strategy are
	positive semidefinite, and the role-indexed blocks multiply according to the
	Kronecker delta on roles.
\end{lemma}
\begin{proof}
	\leanok
	The assertions are the block-diagonal matrix calculations for the direct-sum
	and role-register operators. Positivity follows by writing each positive
	operator as $C^*C$ and placing the corresponding square root in the same
	block decomposition.
\end{proof}

\begin{lemma}[Heterogeneous role-register unsymmetrization]\label{lem:heterogeneous-role-register-unsymmetrization}
	\lean{MIPStarRE.LDT.ProjStrat.qBipartiteConsDefect_roleRegisterProjMeas_arbitrary_eq_average}
	\lean{MIPStarRE.LDT.ProjStrat.qBipartiteConsDefect_extractRoleRegisterBob_le_two_symm}
	\lean{MIPStarRE.LDT.ProjStrat.qBipartiteConsDefect_extractRoleRegisterAlice_le_two_symm}
	\leanok
	\uses{def:projective-strategy, lem:bipartite-strategy-block-identities}
	Let \(\psi\) be a normalized bipartite state on
	\(\mathcal H_{\mathrm A}\otimes\mathcal H_{\mathrm B}\), let
	\(M^{\mathrm A}\) and \(M^{\mathrm B}\) be projective measurements on the two
	prover spaces with a common outcome type, and let \(G\) be an arbitrary
	measurement on the heterogeneous role-register space
	\[
		\{\mathrm A,\mathrm B\}\times(\mathcal H_{\mathrm A}\oplus\mathcal H_{\mathrm B}).
	\]
	After applying the direct-sum role-register symmetrization, the consistency
	defect between the symmetrized point measurement and \(G\) is the average of
	the two defects obtained from the occupied principal blocks of \(G\).  In
	particular, each extracted defect is at most twice the symmetrized defect.
\end{lemma}
\begin{proof}
	\leanok
	The proof is a trace-compression calculation.  The \(AB\)-supported component
	of the symmetrized density pairs with an arbitrary role-register observable
	only through the principal block indexed by
	\((\mathrm A,\mathrm{inl})\) on the left and
	\((\mathrm B,\mathrm{inr})\) on the right, and similarly for the \(BA\)
	component.  Summing over outcomes gives the average-of-two-defects identity,
	and nonnegativity of consistency defects gives the two factor-two
	consequences.
\end{proof}

\begin{lemma}[Heterogeneous role-register symmetrization]\label{lem:heterogeneous-role-register-symmetrization}
	\lean{MIPStarRE.LDT.ProjStrat.roleRegisterSymmStrategy}
	\lean{MIPStarRE.LDT.ProjStrat.roleRegisterSymmStrategy_is_good_three_mul}
	\leanok
	\uses{def:projective-strategy, def:symmetric-projective-strategy, def:good-strategy}
	Let
	\[
		(\psi,A^{\mathrm A},B^{\mathrm A},L^{\mathrm A},
		A^{\mathrm B},B^{\mathrm B},L^{\mathrm B})
	\]
	be a two-space projective strategy passing the low individual degree test
	with error at most \(\eps\).  The heterogeneous role-register construction
	on
	\[
		\{\mathrm A,\mathrm B\}\times
		(\mathcal H_{\mathrm A}\oplus\mathcal H_{\mathrm B})
	\]
	produces a symmetric strategy that is
	\((3\eps,3\eps,3\eps)\)-good.
\end{lemma}
\begin{proof}
	\leanok
	The role-register state is the direct-sum version of the symmetrized state in
	the paper proof.  The axis-parallel and diagonal branches are exactly the
	role averages of the original strategy, and the self-consistency branch is
	the original point-agreement branch.  Since the test chooses among its three
	branches with equal probability, each branch is bounded by \(3\eps\).
\end{proof}

\begin{remark}\label{rem:strategy-convention}
	Throughout this work, the term strategy refers to a projective strategy, and the term symmetric strategy refers to a symmetric projective strategy.
\end{remark}

\begin{definition}[Good strategy]\label{def:good-strategy}
	\lean{MIPStarRE.LDT.SymStrat.IsGood}
	\leanok
	\uses{def:low-individual-degree-test}
	A strategy is $(\eps,\delta,\gamma)$-good if it passes the axis-parallel lines test with probability at least $1-\eps$, the self-consistency test with probability at least $1-\delta$, and the diagonal lines test with probability at least $1-\gamma$.
\end{definition}

\begin{remark}[Good-strategy characterization]\label{rem:good-strat-characterization}
	Using notation which will be introduced in Chapter~\ref{chap:preliminaries} below, a symmetric strategy is $(\eps,\delta,\gamma)$-good if and only if it satisfies the following three conditions. For $\ell$ and $u$ as in the axis-parallel lines test,
	\[
		A_a^u \ot I \simeq_\eps I \ot B^\ell_{[f(u)=a]},
		\quad
		\text{and}
		\quad
		A_a^u \ot I \simeq_\delta I \ot A_a^u.
	\]
	And for $\ell$ and $u$ as in the diagonal lines test,
	\[
		A_a^u \ot I \simeq_\gamma I \ot L^\ell_{[f(u)=a]}.
	\]
\end{remark}

\begin{definition}[Last-direction line notation]\label{not:conditioned-on-last-direction}
	\lean{MIPStarRE.LDT.lastDirectionLine, MIPStarRE.LDT.lastDirectionMeasurementFamily}
	\leanok
	\uses{def:symmetric-projective-strategy}
	For $u \in \F_q^m$, the last-direction axis-parallel line in $\F_q^{m+1}$ is
	\[
		\ell_u = \{(u,x) \mid x \in \F_q\}.
	\]
	If $(\psi,A,B,L)$ is a symmetric strategy for the $(m+1,q,d)$ low individual degree test, we write $B_f^u$ for the operator $B_f^{\ell_u}$. For a function $f \colon \ell_u \to \F_q$, we also write $f(x)$ for $f(u,x)$.
\end{definition}

\begin{definition}[Restricted diagonal lines test]\label{def:restricted-diagonal-lines-test}
	\lean{MIPStarRE.LDT.RestrictedDiagonalSample, MIPStarRE.LDT.MainInductionStep.RestrictedSymStrat.diagonalFailureProbability}
	\leanok
	Consider the $(m,q,d)$ low individual degree test. For $j \in \{1,\dots,m\}$, the $j$-restricted diagonal lines test is the diagonal lines test conditioned on $i=j$. In particular, the $m$-restricted diagonal lines test is the diagonal lines test in which the sampled line is a uniformly random line in $\F_q^m$.
\end{definition}

\section{Formal soundness target}

\begin{theorem}[Quantum soundness of the low individual degree test]\label{thm:main-formal}
	\lean{MIPStarRE.LDT.Test.mainFormal}
	\leanok
	\uses{def:low-individual-degree-test, def:projective-strategy, def:polymeasurements, def:simeq}
	Let $(\psi, A^{\mathrm A}, B^{\mathrm A}, L^{\mathrm A},
	A^{\mathrm B}, B^{\mathrm B}, L^{\mathrm B})$ be a projective strategy that
	passes the $(m,q,d)$ low individual degree test with probability at least
	$1-\eps$. Let \(k>0\) be an integer with \(k \ge 400md\), and set
	\[
		\nu = 100000 k^2 m^4\left(\eps^{1/40000} + (d/q)^{1/40000} + e^{-k/(2560000m^2)}\right).
	\]
	Then there exist projective measurements
	$G^{\mathrm A}, G^{\mathrm B} \in \polymeas{m}{q}{d}$ with the following
	properties:
	\begin{enumerate}
		\item on average over $u \sim \F_q^m$,
		      \[
			      A_a^{\mathrm A,u} \ot I \simeq_\nu I \ot G^{\mathrm B}_{[g(u)=a]},
			      \qquad
			      I \ot A_a^{\mathrm B,u} \simeq_\nu G^{\mathrm A}_{[g(u)=a]} \ot I;
		      \]
		\item
		      \[
			      G_g^{\mathrm A} \ot I \simeq_\nu I \ot G_g^{\mathrm B}.
		      \]
	\end{enumerate}
\end{theorem}
\begin{proof}
	\leanok
	\uses{prop:main-formal-source-reduction, prop:main-formal-source-small-error}
	The saturated-error branch is formalized by
	\texttt{mainFormal\_trivial\_witness}.  The remaining branch is the
	small-error source-boundary proposition
	Proposition~\ref{prop:main-formal-source-small-error}.  The
	factor \(400\) and the nonzero sampling condition are the corrected source
	boundaries recorded in \cite{gap:issue-906-main-formal-k-bound} and
	\cite{gap:issue-422-main-formal-zero-k-boundary}.
\end{proof}

\begin{proposition}[Explicit Lean baseline for the final theorem]
\label{prop:main-formal-explicit-baseline}
	\lean{MIPStarRE.LDT.Test.main_formal_explicit_baseline}
	\leanok
	\uses{thm:main-formal}
	Under exactly the hypotheses of Theorem~\ref{thm:main-formal}, the same two
	projective polynomial measurements satisfy all three conclusions with the
	explicit error
	\[
		\nu_{\mathrm{base}}=
		100000 k^2m^4\left(
			\eps^{1/40000}+(d/q)^{1/40000}
			+e^{-k/(2560000m^2)}\right).
	\]
	This proposition is a Lean-only explicit restatement of the corrected
	source theorem, not a second source claim.
\end{proposition}
\begin{proof}
	\leanok
	\uses{thm:main-formal}
	The assertion follows from Theorem~\ref{thm:main-formal} with its displayed
	value of $\nu$.
\end{proof}

\begin{definition}[Formalization support: linear-triangle final errors]
\label{def:main-formal-linear-triangle-errors}
	\lean{MIPStarRE.LDT.Test.mainFormalLinearTriangleRawError}
	\lean{MIPStarRE.LDT.Test.mainFormalLinearTriangleError}
	\leanok
	For $r,R>0$, write
	\[
		E_{r,R}=\eps^{1/r}+(d/q)^{1/r}+e^{-k/(Rm^2)}.
	\]
	Define the uncapped and capped linear-triangle errors by
	\[
		T=10000k^{1/4}m^{1/2}E_{8192,640000},
		\qquad \nu_{\mathrm{lin}}=\min\{1,T\}.
	\]
	These are formalization-only error functions for the alternative quantitative
	proof of Theorem~\ref{thm:main-formal}.
\end{definition}

\begin{lemma}[Formalization support: nonnegativity of the uncapped error]
\label{lem:main-formal-linear-triangle-error-nonnegative}
	\lean{MIPStarRE.LDT.Test.mainFormalLinearTriangleRawError_nonneg}
	\leanok
	\uses{def:main-formal-linear-triangle-errors}
	If $\eps\ge0$, then $T\ge0$.
\end{lemma}
\begin{proof}
	\leanok
	Every factor and every summand in the definition of $T$ is nonnegative.
\end{proof}

\begin{lemma}[Formalization support: scalar regime in the small-error branch]
\label{lem:main-formal-linear-triangle-small-error-domain}
	\lean{MIPStarRE.LDT.Test.cascade_hypotheses_of_mainFormalLinearTriangleRawError_lt_one}
	\leanok
	\uses{def:main-formal-linear-triangle-errors}
	Assume $\eps\ge0$, $k>0$, and $T<1$.  Then
	\[
		k\ge1,\qquad m\ge1,\qquad 0\le\eps\le1,
		\qquad d\le q,\qquad q>0.
	\]
\end{lemma}
\begin{proof}
	\leanok
	The positivity of $k,m,q$ gives the three boundary inequalities.  If
	$\eps>1$ or $d>q$, the corresponding summand of $E_{8192,640000}$ is at
	least one, while the prefactor in $T$ is at least $10000$, contradicting
	$T<1$.
\end{proof}

\begin{lemma}[Formalization support: three square-root envelope steps]
\label{lem:main-formal-linear-triangle-envelope-eighth-root}
	\lean{MIPStarRE.LDT.Test.stepEnvelope1024_rpow_eighth_le}
	\leanok
	\uses{def:main-formal-linear-triangle-errors}
	Assume $q>0$, $k\ge1$, $m\ge1$, $0\le\eps\le1$, and $d\le q$.  Then
	\[
		E_{1024,80000}^{1/8}\le E_{8192,640000}.
	\]
\end{lemma}
\begin{proof}
	\leanok
	Apply the square-root envelope estimate successively with denominator pairs
	$(1024,80000)$, $(2048,160000)$, and $(4096,320000)$.
\end{proof}

\begin{lemma}[Formalization support: scaled eighth-root estimate]
\label{lem:main-formal-linear-triangle-scaled-eighth-root}
	\lean{MIPStarRE.LDT.Test.rpow_linear_triangle_scale_eighth_le}
	\leanok
	\uses{def:main-formal-linear-triangle-errors}
	Under the same scalar hypotheses, suppose $z\ge0$ and
	\[
		z\le120010k^2m^4E_{1024,80000}.
	\]
	Then
	\[
		z^{1/8}\le\frac92 k^{1/4}m^{1/2}E_{8192,640000}.
	\]
\end{lemma}
\begin{proof}
	\leanok
	\uses{lem:main-formal-linear-triangle-envelope-eighth-root}
	Take eighth roots, use $120010\le(9/2)^8$, and apply the preceding envelope
	estimate.
\end{proof}

\begin{proposition}[Formalization support: complete polynomial measurements]
\label{prop:main-formal-linear-triangle-complete-polynomial-construction}
	\lean{MIPStarRE.LDT.ProjStrat.source_role_register_complete_polynomial_self_consistency_linear_triangle}
	\leanok
	\uses{def:low-individual-degree-test, def:projective-strategy,
		def:polymeasurements, def:simeq}
	Let the projective strategy pass the low individual degree test with failure
	at most $\eps$, and assume $k\ge400md$.  Set
	\[
		s=2I,\qquad z=6s+9\eps+md/q,
	\]
	where $I$ is the main-induction error at $(3\eps,3\eps,3\eps)$.  Then there
	are complete polynomial measurements $G^{\mathrm A},G^{\mathrm B}$ such that
	the two point-evaluation consistencies have error $s$ and their global
	polynomial consistency has error $z$.
\end{proposition}
\begin{proof}
	\leanok
	\uses{thm:main-induction,
		lem:heterogeneous-role-register-unsymmetrization,
		lem:complete-measurement-consistency-triangle}
	Unsymmetrize the polynomial measurement supplied by the main-induction
	theorem.  Apply the complete-measurement linear triangle to the two resulting
	point-consistency estimates and the players' point agreement, then apply the
	Schwartz--Zippel estimate to pass from evaluated to global polynomial
	consistency.
\end{proof}

\begin{proposition}[Formalization support: final projective polynomial measurements]
\label{prop:main-formal-linear-triangle-projective-construction}
	\lean{MIPStarRE.LDT.ProjStrat.source_role_register_final_point_consistency_linear_triangle}
	\leanok
	\uses{def:low-individual-degree-test, def:projective-strategy,
		def:polymeasurements, def:simeq}
	Under the hypotheses of the preceding proposition, let
	\[
		a=100z^{1/4},
		\qquad c=2a+4\sqrt a+2z,
		\qquad \eta=z+\sqrt a,
		\qquad v=6z+6c.
	\]
	There are projective polynomial measurements $Q^{\mathrm A},Q^{\mathrm B}$
	for which both final point-evaluation consistencies have error
	$3(s+\eta+v/2)$, while evaluated and global polynomial consistency have
	error $v/2$.
\end{proposition}
\begin{proof}
	\leanok
	\uses{prop:main-formal-linear-triangle-complete-polynomial-construction,
		lem:orthonormalization-main-lemma,
		prop:completing-to-measurement,
		lem:complete-measurement-consistency-triangle}
	Apply orthogonalization to the two complete polynomial measurements and
	complete the resulting projective submeasurements.  The consistency estimate
	after completion gives the two $\eta$ terms and the $v/2$ polynomial term.
	Two applications of the complete-measurement linear triangle give the final
	point errors.
\end{proof}

\begin{lemma}[Formalization support: absorption of the construction errors]
\label{lem:main-formal-linear-triangle-source-error-absorption}
	\lean{MIPStarRE.LDT.Test.linear_triangle_source_errors_le_uncapped_error}
	\leanok
	\uses{def:main-formal-linear-triangle-errors}
	Assume $\eps\ge0$, $k>0$, and $T<1$, and define $s,z,c,\eta,v$ as above.
	Then
	\[
		3(s+\eta+v/2)\le T,
		\qquad v/2\le T.
	\]
\end{lemma}
\begin{proof}
	\leanok
	\uses{lem:main-formal-linear-triangle-small-error-domain,
		lem:main-formal-linear-triangle-scaled-eighth-root}
	The small-error hypothesis gives $\eps\le1$ and $d\le q$.  The
	main-induction estimate yields
	\[
		z\le120010k^2m^4E_{1024,80000}.
	\]
	The scaled eighth-root estimate and $2221(9/2)\le10000$ imply
	$2221z^{1/8}\le T<1$, and hence $z\le1$.  The orthogonalization and
	completion estimates then bound both displayed construction errors by
	$2221z^{1/8}$.
\end{proof}

\begin{proposition}[Formalization support: saturated linear-triangle branch]
\label{prop:main-formal-linear-triangle-saturated}
	\lean{MIPStarRE.LDT.Test.main_formal_linear_triangle_trivial_witness}
	\leanok
	\uses{def:main-formal-linear-triangle-errors, def:projective-strategy,
		def:polymeasurements, def:simeq}
	For a projective strategy on finite local Hilbert spaces, if $T\ge1$, then
	there are projective polynomial measurements satisfying all three final
	consistency conclusions with error $\nu_{\mathrm{lin}}$.
\end{proposition}
\begin{proof}
	\leanok
	Since $\nu_{\mathrm{lin}}=1$ in this branch, choose arbitrary projective
	polynomial measurements and use the universal upper bound one for normalized
	bipartite consistency defects.
\end{proof}

\begin{proposition}[Formalization support: small-error linear-triangle branch]
\label{prop:main-formal-linear-triangle-small-error}
	\lean{MIPStarRE.LDT.Test.main_formal_linear_triangle_small_error_conclusion}
	\leanok
	\uses{def:main-formal-linear-triangle-errors, def:low-individual-degree-test,
		def:projective-strategy, def:polymeasurements, def:simeq}
	Let the projective strategy pass the low individual degree test with failure
	at most $\eps$.  If $k\ge400md$, $k>0$, and $T<1$, then there are projective
	polynomial measurements satisfying all three final consistency conclusions
	with error $\nu_{\mathrm{lin}}$.
\end{proposition}
\begin{proof}
	\leanok
	\uses{prop:main-formal-linear-triangle-projective-construction,
		lem:main-formal-linear-triangle-source-error-absorption}
	Use the projective measurements from the construction proposition and weaken
	their two explicit errors using the scalar absorption lemma.  Since $T<1$,
	the capped error equals $T$.
\end{proof}

\begin{theorem}[Formalization support: final theorem with the named capped error]
\label{thm:main-formal-linear-triangle-named-error}
	\lean{MIPStarRE.LDT.Test.main_formal_linear_triangle}
	\leanok
	\uses{def:main-formal-linear-triangle-errors, def:low-individual-degree-test,
		def:projective-strategy, def:polymeasurements, def:simeq}
	Under exactly the corrected hypotheses of Theorem~\ref{thm:main-formal},
	there are projective polynomial measurements satisfying its three
	consistency conclusions with error $\nu_{\mathrm{lin}}$.
\end{theorem}
\begin{proof}
	\leanok
	\uses{prop:main-formal-linear-triangle-saturated,
		prop:main-formal-linear-triangle-small-error}
	Split according to whether $T\ge1$.  Apply the saturated proposition in the
	first case and the small-error proposition in the second.
\end{proof}

\begin{theorem}[Complete-measurement linear-triangle bound]
\label{thm:main-formal-linear-triangle}
	\lean{MIPStarRE.LDT.Test.main_formal_linear_triangle_bound}
	\leanok
	\uses{def:low-individual-degree-test, def:projective-strategy,
		def:polymeasurements, def:simeq}
	Let $(\psi,A^{\mathrm A},B^{\mathrm A},L^{\mathrm A},
	A^{\mathrm B},B^{\mathrm B},L^{\mathrm B})$ be a projective strategy that
	passes the $(m,q,d)$ low individual degree test with probability at least
	$1-\eps$.  Let $k>0$ be an integer with $k\ge 400md$, and set
	\[
		T=10000 k^{1/4}m^{1/2}\left(
			\eps^{1/8192}+(d/q)^{1/8192}
			+e^{-k/(640000m^2)}\right),
		\qquad \nu_{\mathrm{lin}}=\min\{1,T\}.
	\]
	Then there are projective measurements
	$Q^{\mathrm A},Q^{\mathrm B}\in\polymeas{m}{q}{d}$ such that, on average
	over $u\sim\F_q^m$,
	\[
		A_a^{\mathrm A,u}\ot I\simeq_{\nu_{\mathrm{lin}}}
			I\ot Q^{\mathrm B}_{[g(u)=a]},
		\qquad
		I\ot A_a^{\mathrm B,u}\simeq_{\nu_{\mathrm{lin}}}
			Q^{\mathrm A}_{[g(u)=a]}\ot I,
	\]
	and
	\[
		Q_g^{\mathrm A}\ot I\simeq_{\nu_{\mathrm{lin}}}
			I\ot Q_g^{\mathrm B}.
	\]
	This is a Lean-only quantitative improvement obtained from the source proof.
	It does not replace or relabel Theorem~\ref{thm:main-formal}.
\end{theorem}
\begin{proof}
	\leanok
	\uses{thm:main-formal-linear-triangle-named-error,
		def:main-formal-linear-triangle-errors}
	Apply Theorem~\ref{thm:main-formal-linear-triangle-named-error} and substitute
	the displayed formulas for $T$ and $\nu_{\mathrm{lin}}$.
\end{proof}

\begin{lemma}[Formalization support: factor-ten comparison]
\label{lem:main-formal-linear-triangle-factor-ten}
	\lean{MIPStarRE.LDT.Test.ten_mul_mainFormalLinearTriangleRawError_le_mainFormalError}
	\leanok
	\uses{def:main-formal-linear-triangle-errors,
		prop:main-formal-explicit-baseline}
	Assume $q>0$, $k\ge1$, $m\ge1$, $0\le\eps\le1$, and $d\le q$.  Then
	\[
		10T\le \nu_{\mathrm{base}}.
	\]
\end{lemma}
\begin{proof}
	\leanok
	Use $8192\le40000$, $640000\le2560000$,
	$k^{1/4}\le k^2$, and $m^{1/2}\le m^4$ to compare the corresponding
	envelopes and prefactors.
\end{proof}

\begin{proposition}[Comparison of the two capped final errors]
\label{prop:main-formal-linear-triangle-comparison}
	\lean{MIPStarRE.LDT.Test.mainFormalLinearTriangleError_le_min_mainFormalError}
	\lean{MIPStarRE.LDT.Test.mainFormalLinearTriangleError_lt_min_mainFormalError}
	\leanok
	\uses{def:main-formal-linear-triangle-errors,
		prop:main-formal-explicit-baseline}
	Assume $\eps\ge0$ and $k>0$.  With $\nu_{\mathrm{lin}}$ as above and
	$\nu_{\mathrm{base}}$ from
	Proposition~\ref{prop:main-formal-explicit-baseline},
	\[
		\nu_{\mathrm{lin}}\le\min\{1,\nu_{\mathrm{base}}\}.
	\]
	If $\nu_{\mathrm{base}}<1$, then the inequality is strict.  More precisely,
	in that branch the formal scalar proof establishes
	$10T\le\nu_{\mathrm{base}}$; the exponential summand makes $T>0$.
\end{proposition}
\begin{proof}
	\leanok
	\uses{lem:main-formal-linear-triangle-factor-ten,
		lem:main-formal-linear-triangle-error-nonnegative}
	When the previous explicit error bound is at least one, both sides are capped
	by one.  Otherwise
	$0\le\eps\le1$ and $d/q\le1$.  Monotonicity of real powers and the exponential
	tail, together with $k^{1/4}\le k^2$ and $m^{1/2}\le m^4$, gives the stated
	factor-ten comparison.
\end{proof}

\begin{proposition}[Source-boundary reduction for the final theorem]\label{prop:main-formal-source-reduction}
	\lean{MIPStarRE.LDT.Test.mainFormalConclusion}
	\leanok
	This proposition records the saturated-error reduction for the corrected
	two-space theorem
	Theorem~\ref{thm:main-formal}.  It proves the saturated-error branch by the
	trivial consistency bound and reduces the remaining non-vacuous branch to
	Proposition~\ref{prop:main-formal-source-small-error}.  It is not
	an additional hypothesis of the paper theorem.
\end{proposition}
\begin{proof}
	\leanok
	\proves{prop:main-formal-source-reduction}
	\uses{prop:main-formal-source-small-error}
	The saturated-error branch is already proved.  The small-error branch is
	Proposition~\ref{prop:main-formal-source-small-error}.
\end{proof}

\begin{proposition}[Small-error source-boundary conclusion for the final theorem]\label{prop:main-formal-source-small-error}
	\lean{MIPStarRE.LDT.Test.mainFormal_smallErrorConclusion}
	\leanok
	This is the non-vacuous branch of the corrected two-space theorem.
	It proves the conclusions of Theorem~\ref{thm:main-formal} under
	\(k\ge 400md\), \(k>0\), and \(\nu<1\).  It is not a source assumption.  Its
	proof uses the two-space role-register reduction and the scalar absorption
	available after excluding the zero-sampling boundary.
\end{proposition}
\begin{proof}
	\leanok
	\proves{prop:main-formal-source-small-error}
	\uses{prop:main-formal-source-two-space-role-register}
	The proof applies the checked two-space role-register construction and then
	absorbs its explicit errors into \(\nu\) at the nonzero sampling boundary.
\end{proof}

\begin{proposition}[Two-space role-register reduction for the source theorem]\label{prop:main-formal-source-two-space-role-register}
	\lean{MIPStarRE.LDT.Test.mainFormalConclusion_ofRoleRegisterScalarBoundary}
	\leanok
	\uses{lem:heterogeneous-role-register-unsymmetrization}
	The printed final theorem starts from a projective strategy on two possibly
	different local Hilbert spaces.  The source proof uses the role-register
	symmetrization and then unsymmetrizes the polynomial measurement.  Lean
	formalizes this reduction for the general two-space source strategy.  The
	trace-level factor-two unsymmetrization estimate for arbitrary heterogeneous
	role-register measurements is formalized in
	Lemma~\ref{lem:heterogeneous-role-register-unsymmetrization}.  The
	heterogeneous role-register symmetrization may be fed into the source-shaped
	main-induction theorem under the corrected large-\(k\) hypothesis
	\(k\ge 400md\), and the resulting role-register consistency estimate can be
	unsymmetrized into the two point-consistency estimates on the original
	two-space strategy with the expected factor \(2\).

	The remaining steps of the source-boundary passage are also checked.  The
	point-agreement branch is available for the paper-faithful two-space
	strategy, and the Schwartz--Zippel Step~5 theorem has been generalized from
	the same-space carrier to a bipartite state on \(H_A\otimes H_B\).  The
	heterogeneous triangle/SDD comparison theorem gives the complete-measurement
	full-polynomial consistency statement at the end of the paper's Step~5
	calculation.  The two heterogeneous orthonormalization applications produce
	projective submeasurements on the two local Hilbert spaces, and the completion
	step widens the resulting left- and right-factor state-dependent-distance
	estimates to the orthonormalize-and-complete error appearing in the paper.
	The repaired polynomial line-169 consistency relations use the
	Cauchy--Schwarz loss from the pre-completion orthonormalization estimates,
	not a new hypothesis.  The final point-evaluation triangle postprocesses
	these relations and the completed \(Q_A,Q_B\) projective consistency estimate
	by evaluation at a point, combines them by the heterogeneous triangle
	inequality, and absorbs the explicit pre-absorption errors into the final
	\(\nu=\texttt{mainFormalError}\) bound under the nonzero scalar-cascade
	boundary \(0<k\).  This is the small-error branch used in the corrected
	two-space source theorem, where \(k>0\) is part of the theorem statement.
\end{proposition}
\begin{proof}
	\proves{prop:main-formal-source-two-space-role-register}
	\lean{MIPStarRE.LDT.Test.mainFormalConclusion_ofRoleRegisterScalarBoundary}
	\leanok
	\uses{lem:heterogeneous-role-register-symmetrization, lem:heterogeneous-role-register-unsymmetrization}
	\uses{thm:main-induction}
	The proof is the linked heterogeneous role-register chain, ending in
	\texttt{mainFormalConclusion\_ofRoleRegisterScalarBoundary}.  It
	combines the source-shaped main-induction theorem, the heterogeneous
	unsymmetrization estimates, the step making measurements projective on both
	local Hilbert spaces, the completion and line-169 consistency steps, and the
	final scalar absorption under the nonzero sampling boundary.
\end{proof}

\begin{remark}[Final-theorem zero-sampling boundary]\label{rem:main-formal-source-k-range-boundary}
	The paper states Theorem~\ref{thm:main-formal} for \(k\ge md\), while the
	corrected formal route uses the confirmed large-\(k\) condition
	\(k\ge 400md\) and a nonzero \(k\) boundary for the scalar cascade.  The
	factor \(400\) is now treated as a statement correction, following
	Theorem~\ref{thm:main-induction}.  The other corrected final-theorem boundary is
	the zero-sampling case allowed when \(d=0\), since \(400md\le k\) then permits
	\(k=0\).  At this excluded boundary the displayed final error itself collapses to \(0\);
	the formal calculation of this equality is linked in Lean, and the
	obstruction is documented in \cite{gap:issue-422-main-formal-zero-k-boundary}.
\end{remark}

\begin{remark}[$k$-bound boundary for the public theorem]\label{rem:main-formal-k-boundary}
	The paper statement uses the weaker hypothesis $k \ge md$, but the proof
	invokes the Section~6 pasting theorem with side condition $k \ge 400\,md$.
	This side condition is not implied by $k \ge md$.  The Lean theorem
	therefore records the stronger large-$k$ assumption required by the pasting
	step, together with the scalar-cascade boundary \(k>0\), rather than hiding
	the side-condition gap.  The helper
	\texttt{mainFormal\_trivial\_witness} remains only for already-saturated
	branches where $\nu \ge 1$ makes the three $\simeq_\nu$ conclusions vacuous;
	it is not a replacement for the zero-sampling boundary \(k=0\) when \(md=0\).
	The source-boundary reduction proposition proves the saturated-error branch
	and delegates the non-vacuous branch to the checked small-error proposition;
	neither declaration is an extra hypothesis of Theorem~\ref{thm:main-formal}.
\end{remark}
